\cleardoublepage
\chapter*{Додатки}
\addcontentsline{toc}{chapter}{Додатки}
\markboth{Додатки}{Додатки}

\section*{Додаток А. Короткий словник ключових термінів}
\addcontentsline{toc}{section}{Додаток А. Короткий словник ключових термінів}

\begin{description}[leftmargin=0mm,itemsep=2pt,font=\normalfont\bfseries]
\item[Злоякісний програмний код (ЗПК).] Програмне забезпечення, створене для несанкціонованих, шкідливих дій у комп'ютерній системі.
\item[Сигнатура.] Формалізований опис ознак, наявність яких свідчить про належність зразка до відомої загрози.
\item[Статичний аналіз.] Дослідження зразка без його виконання.
\item[Динамічний аналіз.] Дослідження зразка під час виконання в контрольованому середовищі.
\item[Пісочниця.] Ізольоване інструментоване середовище для виконання й спостереження зразка.
\item[Обфускація.] Перетворення коду, що зберігає функцію, але ускладнює розуміння.
\item[Пакування.] Перетворення виконуваного файлу на стиснений/зашифрований образ із розпакувальником.
\item[Поліморфізм.] Відтворення коду в різних байтових формах зі збереженням функції.
\item[Індикатор компрометації (IoC).] Спостережувана ознака ймовірної компрометації.
\item[Розвідка загроз.] Засноване на доказах знання про загрози, придатне для прийняття рішень.
\item[APT.] Цілеспрямована стала загроза --- ресурсоване, наполегливе, адаптивне угруповання.
\item[Полювання за загрозами.] Проактивний, керований гіпотезами пошук прихованих загроз.
\item[Час присутності (dwell time).] Проміжок від компрометації до виявлення.
\item[Кіберстійкість.] Здатність витримувати атаки, обмежувати вплив і швидко відновлюватися.
\item[Ешелонований захист.] Поєднання кількох незалежних рубежів, що компенсують слабкості один одного.
\end{description}

\section*{Додаток Б. Зіставлення методів виявлення}
\addcontentsline{toc}{section}{Додаток Б. Зіставлення методів виявлення}

\begin{center}\footnotesize
\begin{tabularx}{\linewidth}{L{26mm} Y Y}
\toprule
Метод & Сильні риси & Обмеження \\
\midrule
Сигнатурне виявлення & точність, швидкість для відомого & безсиле проти нового, пакування \\
Евристичний аналіз & виявляє невідоме & хибні тривоги \\
Поведінкове виявлення & інваріантне до запису коду & потребує виконання, телеметрії \\
Машинне навчання & узагальнення, масштаб & дані, дрейф, змагальність \\
Виявлення за подібністю & працює із сімействами & вразливе до обфускації \\
Полювання за загрозами & знаходить приховане, APT & трудомісткість, кваліфікація \\
\bottomrule
\end{tabularx}
\end{center}

\section*{Додаток В. Контрольний перелік зрілості протидії ЗПК}
\addcontentsline{toc}{section}{Додаток В. Контрольний перелік зрілості протидії ЗПК}

\begin{enumerate}[leftmargin=6mm,itemsep=2pt]
\item Визначено модель загроз і пріоритетні активи (розділ~\ref{chap:ch04}).
\item Забезпечено профілактику: гартування, контроль застосунків, найменші привілеї, оновлення (розділ~\ref{chap:ch10}).
\item Розгорнуто багаторівневе виявлення: сигнатури, поведінка, машинне навчання (розділи~\ref{chap:ch06}--\ref{chap:ch11}).
\item Налагоджено централізовану телеметрію кінцевих точок і мережі (розділ~\ref{chap:ch10}).
\item Інтегровано розвідку загроз і забезпечено обмін індикаторами (розділ~\ref{chap:ch12}).
\item Впроваджено проактивне полювання за загрозами (розділ~\ref{chap:ch13}).
\item Забезпечено цілісність ланцюга постачання й прозорість складу ПЗ (розділ~\ref{chap:ch15}).
\item Підготовлено план реагування на інциденти й форензичного збору (розділ~\ref{chap:ch16}).
\item Сформовано культуру безпеки й навчання персоналу (розділ~\ref{chap:ch17}).
\item Запроваджено метрики, модель зрілості й цикл безперервного вдосконалення (розділ~\ref{chap:ch17}).
\end{enumerate}
